\documentclass[12pt,a4paper]{article}
\usepackage[utf8]{inputenc}
\usepackage[T1]{fontenc}
\usepackage{enumitem}
\usepackage[margin=2.5cm]{geometry}
\usepackage{parskip}
\usepackage{amsmath}
\usepackage{xcolor}

\title{plsssssssss — Final Exam --- Answer Key}
\date{}

\begin{document}

\maketitle

{\large \textbf{\textcolor{red}{ANSWER KEY --- Not for Distribution}}}

\vspace{0.5em}



\noindent
\textbf{Total points:} 6
\hfill
\textbf{Number of questions:} 6

\vspace{1em}
\hrule
\vspace{1.5em}


\noindent
\textbf{Question 1 (1.00 pts).}~
Which of the following is equivalent to the statement that a neural network is a universal approximator?


\begin{enumerate}[label=\textbf{\Alph*.}]

  \item It can only approximate linear functions.

  \item \textbf{\textcolor{teal}{[CORRECT]}}  It can approximate any continuous function.

  \item It can only approximate discrete functions.

  \item It can only approximate polynomial functions.

\end{enumerate}



\textit{\small Explanation: A universal approximator can approximate any continuous function, which is a key property of neural networks.}



\vspace{0.8em}
\noindent\rule{\linewidth}{0.2pt}
\vspace{0.3em}


\noindent
\textbf{Question 2 (1.00 pts).}~
Which of the following is a consequence of the GIGO principle?


\begin{enumerate}[label=\textbf{\Alph*.}]

  \item The output will always be accurate.

  \item The input must be numeric.

  \item \textbf{\textcolor{teal}{[CORRECT]}}  The output will not make sense if the input is incorrect.

  \item The network will always converge.

\end{enumerate}



\textit{\small Explanation: The GIGO principle states that 'Garbage In, Garbage Out', meaning that if the input to a computer does not make sense, the output will also not make sense.}



\vspace{0.8em}
\noindent\rule{\linewidth}{0.2pt}
\vspace{0.3em}


\noindent
\textbf{Question 3 (1.00 pts).}~
Each technology in the relationship among technologies is unique and interrelated.


\textbf{Answer:} \textcolor{teal}{True}



\textit{\small Explanation: The context states that each technology is unique and interrelated.}



\vspace{0.8em}
\noindent\rule{\linewidth}{0.2pt}
\vspace{0.3em}


\noindent
\textbf{Question 4 (1.00 pts).}~
A neural network has been trained to make decisions based on the opinions of expert underwriters. What is the result of using this neural network?


\textbf{Model Answer:} \textcolor{teal}{A 12\% reduction in delinquencies compared with human experts.}



\textit{\small Explanation: neural network; expert underwriters; 12\% reduction}



\vspace{0.8em}
\noindent\rule{\linewidth}{0.2pt}
\vspace{0.3em}


\noindent
\textbf{Question 5 (1.00 pts).}~
How do artificial neural networks provide human-comprehensible rules and insights into complex processes?


\textbf{Model Answer:} \textcolor{teal}{Artificial neural networks provide human-comprehensible rules and insights into complex processes through knowledge extraction, which can yield governing rules and insights into complex processes.}



\textit{\small Explanation: Knowledge extraction; Governing rules; Human-comprehensible insights}



\vspace{0.8em}
\noindent\rule{\linewidth}{0.2pt}
\vspace{0.3em}


\noindent
\textbf{Question 6 (1.00 pts).}~
Explain how changes in Total Phosphorus (TP) affect water transparency, and discuss the implications of this relationship on eutrophication.


\textbf{Key Points:} \textcolor{teal}{I. Introduction to the relationship between TP and water transparency, II. Analysis of the impact of TP on Secchi Disk Depth (SDD) and Total Suspended Solids (TSS), III. Discussion of the implications of this relationship on eutrophication.}



\textit{\small Explanation: Address the relationships between TP, Secchi Disk Depth (SDD), and Total Suspended Solids (TSS) at a constant CHL.}





\end{document}
